% Generated draft: requires mathematical and source review.
\documentclass[11pt]{article}
\usepackage[margin=25mm]{geometry}
\usepackage{fontspec}
\setmainfont{Latin Modern Roman}
\usepackage{amsmath,amssymb,mathtools,unicode-math}
\setmathfont{Latin Modern Math}
\usepackage{enumitem,needspace,xurl,hyperref}
\setlength{\parindent}{0pt}
\setlength{\parskip}{5pt}
\setlength{\emergencystretch}{3em}
\newcommand{\N}{\mathbb{N}}
\newcommand{\Z}{\mathbb{Z}}
\newcommand{\Q}{\mathbb{Q}}
\newcommand{\R}{\mathbb{R}}
\newcommand{\C}{\mathbb{C}}
\newcommand{\F}{\mathbb{F}}
\newcommand{\A}{\mathbb{A}}
\newcommand{\PP}{\mathbb{P}}
\newcommand{\E}{\mathbb{E}}
\newcommand{\eps}{\varepsilon}
\newcommand{\dd}{\,\mathrm{d}}
\begin{document}
\section*{UnsolvedMath ID: 600009}
\textbf{Source problem code:} catalogue entry
\subsection{Definitions and mathematical statement}
Consider the origami hyperbolic-paraboloid crease pattern referred to in Section 8 of Tabachnikov's source. In addition to its displayed creases, insert an ``invisible'' straight fold along a chosen diagonal of every elementary trapezoid. For a specified crease pattern, diagonal choices, and folding configuration, let $S$ denote the resulting piecewise-linear surface, with each face planar.
\begin{enumerate}
\item Determine the shape of $S$ as the pattern folds. The source asks for the shape of the folded surface, rather than explicitly only for a refinement limit.
\item Determine the surfaces produced by analogous constructions from other patterns of fold lines; the supplied research summary mentions circular pleats as an example.
\end{enumerate}
The supplied statement does not give the crease coordinates, boundary, compatible folding constraints, parameterization of folding depth, or a mathematical meaning of ``invisible.'' Thus it does not determine a unique surface from the text alone. The second question also leaves the class of alternative patterns and the intended notion of determining a shape unspecified.

The supplied later research summary reports a result for a more specified triangulated construction with one added diagonal per trapezoid: it is generically a one-degree-of-freedom origami, and as the pleats refine its folded shapes converge, at each folding depth, to a hyperbolic paraboloid of the form
\[
z=k(x^2-y^2),
\]
where $k$ is a monotone function of the folding angle. The summary states independence of square spacing and of diagonal choices \emph{among the symmetric schemes}. It supplies neither the exact admissible schemes nor a convergence topology or angle-to-$k$ formula. This limiting claim does not itself identify a finite piecewise-linear surface with the smooth quadratic graph. The summary also reports bistability with two mirror-symmetric stable states in the cited work, without supplying its mechanical stability model. It regards the first question as resolved but the second as only partially addressed by $N$-gon hypar extensions and framework remarks.
\subsection{Short English statement}
The problem asks for the shapes obtained by folding a diagonal-triangulated origami hypar pattern and analogous patterns. The supplied 2019 research summary reports a hyperbolic-paraboloid refinement limit for specified symmetric constructions, while the broader question about other crease patterns remains open. The precise finite folding model and limiting formulation need the original sources.
\subsection{Sources}
\begin{itemize}
\item Local UnsolvedMath registry, permanent ID 600009, exact source code catalogue entry, ``Baker's Dozen --- Origami hyperbolic paraboloids.'' The registry labels it solved and SOLVED-IN-LITERATURE, while its research summary expressly leaves the second question open.
\item Tabachnikov, \emph{A Baker's Dozen of Problems} (2015), Section 8, unnumbered questions, \url{https://amj.math.stonybrook.edu/html-articles/Files-2015-2024/14-01/}. The record says this was collected from an AMR-indexed open-problem list by publisher-HTML extraction, accessed July 29, 2026, and that the source presented it as open.
\item Liu--Tachi--Paulino, \emph{Nature Communications} 10:4238 (2019), cited in the supplied research summary for the generic one-degree-of-freedom triangulated origami, the hyperbolic-paraboloid limit and its stated independence properties, and bistability. A paper title, DOI, and precise theorem hypotheses are not supplied; no further bibliographic details are invented here.
\end{itemize}
\end{document}
